\documentclass[aspectratio=169]{beamer}
\usepackage[utf8]{inputenc}
\usepackage[T1]{fontenc}
\usepackage[brazil]{babel}
\usepackage{amsmath}
\usepackage{booktabs}
\usepackage{listings}
\usepackage{graphicx}

\usetheme{Madrid}
\usecolortheme{default}

\title[Lynceus]{Lynceus}
\subtitle{Motor de Extração eBPF/XDP para Ambientes DDoS}
\author{Leonardo Herkenhoff}
\institute[PPComp/IFES]{Mestrado Profissional em Computação Aplicada (PPComp) -- IFES}
\date{\today}

\begin{document}

\frame{\titlepage}

\begin{frame}{O Problema}
    \textbf{Contexto e Desafio}
    \begin{itemize}
        \item \textbf{Limitações Atuais:} Extratores clássicos (NetFlow/sFlow) operam no espaço de usuário (\textit{user-space}) ou via amostragem (\textit{sampling}), omitindo a granularidade microscópica necessária para Aprendizado de Máquina.
        \item \textbf{Problema Central:} Centros de mitigação DDoS exigem captura de pacotes na velocidade máxima da interface (\textit{wire-speed}), sem esgotar I/O ou memória do servidor.
        \item \textbf{Proposta (Lynceus):} Arquitetura híbrida. O \textit{kernel} (via eBPF/XDP) disseca os pacotes e o \textit{user-space} consolida a taxonomia estatística, resultando em \textbf{495 atributos} (\textit{features}) extraídos em tempo real.
    \end{itemize}
\end{frame}

\begin{frame}{Arquitetura de Extração}
    O processamento é dividido em dois planos concorrentes:
    \vspace{0.3cm}
    
    \textbf{1. Plano de Dados (\textit{Data Plane} - Kernel-Space)}
    \begin{itemize}
        \item Análise de pacotes L2 a L7 diretamente no buffer do \textit{driver} da placa de rede.
        \item \textbf{Desencapsulamento Nativo:} Extração do conteúdo interno de túneis GRE e VXLAN.
        \item \textbf{Descoberta L7:} Acesso à carga útil (\textit{payload}) para identificar vetores volumétricos críticos (DNS, SNMP, NTP, SSDP).
    \end{itemize}
    \vspace{0.2cm}
    
    \textbf{2. Plano de Controle (\textit{Control Plane} - User-Space)}
    \begin{itemize}
        \item Processamento assíncrono e massivo dos pacotes classificados pelo \textit{kernel}.
        \item \textbf{Exportação Agnóstica:} Envio de dados diretamente via \texttt{stdout}, neutralizando gargalos de escrita em disco.
    \end{itemize}
\end{frame}

\begin{frame}{Superando Restrições do Kernel}
    O verificador do eBPF (\textit{Verifier}) bloqueia execuções lentas, loops ilimitados e complexidade de memória $O(N)$.
    \vspace{0.3cm}
    
    \textbf{Solução 1: Verificação de Limite Único (\textit{Single Boundary Check})}
    \begin{itemize}
        \item A dissecação restringe-se aos primeiros 64 bytes de cada pacote. Isso fixa o custo computacional em $O(1)$, independente do tamanho real do tráfego.
    \end{itemize}

    \textbf{Solução 2: Independência Estrutural (\textit{Purity})}
    \begin{itemize}
        \item Para assegurar portabilidade em kernels LTS (5.15+) sem falhas de biblioteca (BTF), funções críticas como \texttt{memcmp} foram escritas nativamente (\textit{unrolled}) no código BPF.
    \end{itemize}
\end{frame}

\begin{frame}{Estatísticas em Tempo Contínuo}
    Cálculos de fluxo não podem armazenar o histórico completo de pacotes devido à limitação de memória RAM.
    
    \vspace{0.3cm}
    \textbf{1. Algoritmo de Welford}
    \begin{itemize}
        \item Atualiza recursivamente a Média ($\mu$), Variância ($\sigma^2$), Assimetria e Curtose a cada novo pacote lido, viabilizando métricas de 4ª ordem.
    \end{itemize}
    
    \vspace{0.2cm}
    \textbf{2. Estimador P2}
    \begin{itemize}
        \item Extrair a Mediana exigiria a ordenação de vetores em $O(N \log N)$. A heurística $P^2$ aproxima percentis em tempo real com uso constante de memória $O(1)$.
    \end{itemize}
\end{frame}

\begin{frame}{Alta Performance e Paralelismo}
    \textbf{O Desafio da Concorrência}
    \begin{itemize}
        \item Centralizar os eventos do \textit{kernel} em um único buffer gera contenção de concorrência (\textit{lock contention}) em arquiteturas \textit{multi-core}.
    \end{itemize}
    
    \textbf{Solução: Canais Independentes}
    \begin{itemize}
        \item A infraestrutura aloca um buffer de comunicação dedicado para cada núcleo lógico da CPU (\texttt{ARRAY\_OF\_MAPS}).
        \item O código eBPF identifica o núcleo atual (\texttt{bpf\_get\_smp\_processor\_id()}) e despacha os dados no canal correspondente.
        \item Processos \textit{workers} em \textit{user-space} efetuam a leitura paralela e livre de bloqueios (\textit{lockless}), aproximando-se da escalabilidade linear.
    \end{itemize}
\end{frame}

\begin{frame}{Validação e Próximos Passos}
    \textbf{O Ciclo de Experimentação (\textit{Extract-Label-Purge})}
    \begin{itemize}
        \item Automação inicializa redes efêmeras virtuais (\texttt{veth pairs}).
        \item Injeção de bancos de dados (\textit{datasets}) de DDoS na velocidade nominal da placa.
        \item O sistema gera as matrizes de características matemáticas, rotula e purga as conexões virtuais.
    \end{itemize}
    
    \textbf{Resultados Finais}
    \begin{itemize}
        \item Entrega estruturada de \textbf{495 atributos de rede} por fluxo analisado.
        \item \textbf{Trabalhos Futuros:} Adição de bloqueio ativo no plano de dados e otimização de exportação usando Memória Compartilhada Interprocessos (SHM IPC).
    \end{itemize}
\end{frame}

\end{document}
